\section*{\centering{ABSTRACT}}

This work investigates the performance of three tree-based \textit{Gradient Boosting} algorithms --- XGBoost, LightGBM, and CatBoost --- in a Federated Learning environment applied to bank account opening fraud detection. Using the BAF (Bank Account Fraud) dataset and the Flower framework for federated simulation, two aggregation strategies --- Bagging and Cyclic --- were evaluated with 3 clients, 5 communication rounds, and balanced IID partitioning. The federated models achieved TPR between 54.83\% and 56.79\% at the 5\% FPR operating point, results competitive with the centralized BAF baseline (40--55\%). XGBoost Bagging achieved the best absolute performance (56.79\% TPR, AUC 0.8934), while CatBoost Cyclic reached a nearly equivalent result (56.72\% TPR) with lower communication cost (1.27 MB). Applying independent thresholds per demographic group raised the FPR Ratio from $\approx$0.3 to $\approx$0.88--1.0, demonstrating that algorithmic fairness can be achieved with an average loss of only 2.0 percentage points in TPR. The results demonstrate the viability of tree-based models for federated fraud detection systems, offering superior communication efficiency compared to deep neural network approaches.

{\bf Keywords:} Federated Learning, Bank Fraud Detection, XGBoost, LightGBM, CatBoost, Flower Framework, Data Privacy, Tree-Based Models.
